\documentclass{article}
\usepackage{graphicx} % Required for inserting images
\usepackage{float}
\usepackage{amsmath}
\title{Lesson 7 Homework (Extra Challenge) - Proof Writing}
\author{William Wang}
\date{October 2025}

\begin{document}

\maketitle

\section{}
We have
\[\arccos(\sin \frac{6\pi}{7}) = \theta \implies \sin(\frac{6\pi}{7}) = \cos \theta\]
Since we know that the graph of $\sin$ is just a transformation of $\cos$, we can actually express this equation under using the same trigonometric function. We have 
\[\sin(x+\frac{\pi}{2}) = \cos x\]
Hence we have 
\[\cos \theta = \sin(\theta + \frac{\pi}{2}) = \sin(\frac{6\pi}{7})\]
So we have 
\[\theta + \frac{\pi}{2} = \frac{6\pi}{7} + 2a\pi\]
Where \(a\) is a constant for the cycle of sine. However, the solution is restricted by the range of \(\arccos x \in [0, \pi] \). Hence the only solution occurs when $a = 0$. So
\[\theta +\frac{\pi}{2} = \frac{6\pi}{7} \implies \theta = \frac{5\pi}{14}\]
Which is in the range of $\arccos$. 

\end{document}
